% Source PDF page 527 (printed p. 513)
\section*{文献说明}
\addcontentsline{toc}{section}{文献说明}
关于无约束最小化与非线性方程的算法，Dennis 和 Schnabel [DS96] 以及 Ortega 和 Rheinboldt [OR00] 是两部标准参考书。在强凸性与 Hessian 的 Lipschitz 连续性假设下得到的二次收敛结论，归于 Kantorovich [Kan52]。Polyak [Pol87，第1.6节] 对涉及未知常数的收敛结论的作用，给出了颇有见地的讨论，第9.5.3节导出的结论就属于这一类。

自协调函数由 Nesterov 和 Nemirovski [NN94] 引入。第9.6节和习题9.14至9.20中的全部结论，都可以在他们的书中找到，不过通常形式更一般，或者记号不同。Renegar [Ren01] 简明而优美地介绍了自协调函数，以及它们在原始—对偶内点算法分析中的作用。Peng、Roos 和 Terlaky [PRT02] 从自正则函数的角度研究内点法，这一类函数与自协调函数相似，但并不相同。

第9.7节相关内容的参考文献，见附录C末尾。
